% Page budget: 1 page
\chapter{סיכום}

\section{ממצאים עיקריים}
הדוח מצא פער ברור במחקר ובשוק: המחקרים העדכניים כבר שומרים על עקביות חזותית ואף מנהלים מצב סיפור בתוך הפקה אחת \cite{huang2026dramaagent}, ומוצרים מסחריים שומרים דמויות ועולם בין פרקים; אולם אף פתרון אינו מנהל זיכרון עלילתי מתמשך של סדרה שפרקיה מופקים בנפרד לאורך זמן, בשליטת היוצר ועם זיהוי סתירות בין פרקים, ואף אחד אינו נגיש ליוצרים בעברית. בבדיקת ארבעת כלי הווידאו בפרק~4 נמצא שהם שומרים על מראה הדמות, אך לפי סקירת דפי המוצר אינם מציגים זיהוי סתירות עלילתיות, ואף אחד מהם לא היה ניתן לשימוש בעברית בתוכנית חינמית. מבין שלוש חלופות התכן, חלופה ג׳ ההיברידית קיבלה בהערכת הצוות ציון משוקלל של \textenglish{3.50}, לעומת \textenglish{3.05} לחלופה ב׳ ו־\textenglish{2.85} לחלופה א׳, והיא ניצחה בכל תרחישי בדיקת הרגישות. תוכנית העבודה כוללת חמש חבילות עבודה, ארבעה עשר שבועות ושתי אבני דרך. ארבעה סיכונים נמצאו קריטיים: \mbox{\textenglish{R1, R2, R3}} ו־\textenglish{R7}.

\section{מוכנות לשלב התכן}
הבחירה בארכיטקטורה, מטריצת ההחלטה, תוכנית העבודה ומטריצת הסיכונים מוכנות לתחילת התכן המפורט, משום שכל אחד מהם נשען על דרישות מוגדרות מפרק~2 ועל מקורות או הערכות צוות המסומנים במפורש. עם זאת, ערכי היעד בפרק~2 הם יעדי תכנון של הצוות, וקו הבסיס של התהליך הקיים יימדד בחבילת עבודה~1. מחיר \textenglish{CapCut} בפרק~4 נלקח ממקור משני. הציון \textenglish{3.50} והאומדנים הכספיים נשארים הערכה של הצוות ולא עובדה שנמדדה.

\section{פעולות מיידיות לאחר האישור}
\begin{enumerate}
  \item למדוד את קו הבסיס של התהליך הידני על שלושה כלים קיימים, בשבועות 1 עד 2.
  \item לאמת את מחיר \textenglish{CapCut} מול דף המחירים הרשמי.
  \item להתחיל בחבילת עבודה~1, ובחבילת עבודה~3 להשוות איכות שוט מקומי מול שוט דרך \textenglish{API}.
\end{enumerate}

\begin{directornote}
המסקנה היא שהפרויקט מוכן לצאת לדרך, ושלושת הצעדים הללו סוגרים את הפערים שנותרו לפני אבן הדרך הראשונה.
\end{directornote}
